\subsection{Finding Percent Change} 
We might also compute the percent change for an exponential model described in some other terms.
\begin{example}
10,000 copies of a limited edition collectible card are initially sold. After 15 years, only 10 cards in good condition remain in circulation. Wear and tear on collectibles follows an exponential decay model. What percent of the cards are destroyed each year? (Round to one decimal place.)
\begin{itemize}
\item In this problem, we must find the multiplier first.
\item The multiplier is \makeblank.
\item The time required to multiply by this is number is \makeblanksmall years.
\item There are \makeblank[1in] cards in circulation after \makeblanksmall years.
\item We can write the number of cards in circulation after $t$ years as $n(t)=\makeblank$
\item We want to know the percent lost \emph{each year}. We rewrite the exponential part of the expression as \makeblank.
\item The multiplier each year is \makeblank.
\item What do we need to subtract from 100\% to get \makeblanksmall\%? \makeblank[1in]
\item \makeblanksmall\% of the cards are destroyed each year.
\end{itemize}
\end{example}
Let's see an example of the work we would typically write out in a problem like this.
\begin{example}
After an economic boom, the population of a small town doubles over the course of 5 years. By what percent did the population increase each year? (Round to one decimal place.)
\begin{lpic}{}{7}
\end{lpic}
\end{example}